% ---------------------------------------------------------------------------
% Experiment History -- subsection 8: Realistic Channel, Detector Suite, and
% Matched Detectability.  Sources: README A.7.
% Jobs: 101870 (m1), 102305 (m1b), 102316 (m2 det), 102319 (m2 suite),
% 102390 (dense re-sweep).
%
% Rewritten 2026-09-11: cut from ~830 words to ~360, to the register of the
% user-written parts 2 and 4.
%
% This part retracts two of our own headline numbers -- "+70% channel-aware"
% and "stealthy BER 0.065-0.11" -- and both have the SAME cause: a comparison
% that held the wrong thing fixed.  Keep that common cause visible; it is the
% point of the subsection and the justification for the two reporting rules
% in the System Model, which are NOT restated here (convention (f)).
%
% Figure: artifacts/sim08/frontier_dense/matched_detectability.png, CROPPED
% TO TWO PANELS (5 dB, 30 dB).  Five is overkill.
% A second candidate, artifacts/sim08/frontier/stealth_suite_vs_snr.png,
% HAS NOT BEEN EYEBALLED -- check before committing to it.
% ---------------------------------------------------------------------------

\subsection{Realistic Channel, Detector Suite, and Matched Detectability}
\label{subsec:exphist:realistic}
The previous step ended on a caveat about its own channel, so we replaced it: an independent
frequency-selective block-fading realisation per link, from a tapped-delay-line model at 100\,ns spread
and 5.2\,GHz, with noise at a target $E_b/N_0$ and zero-forcing equalisation of the legitimate link.
The jammer now reaches the receiver through its own channel, so after equalisation its effect on
subcarrier $k$ scales with $G_j[k]/H_0[k]$ and is amplified wherever the legitimate link is faded.

The chain behaves, with clean BER waterfalling from $0.088$ at 5\,dB to $0.0003$ at 30\,dB.
A sparse in-band jammer then imposes a BER floor that does not move with SNR, roughly $0.05$--$0.07$ at
eight active subcarriers and $0.34$ broadband, across the whole range.
That is the strongest result here and the only one needing no detector and no stealth threshold: the
impairment is interference, not noise, so the link cannot buy its way out with transmit power, and the
attacker is relatively strongest exactly where the link is cleanest.
We also measured channel-aware selection beating blind selection by up to $70\,\%$ at equal power, which
is retracted below.

Folding the energy detector back in, calibrated per SNR on faded clean frames, recovers what the
lossless channel had ruled out: an attacker holds $\Pr[\text{flag}] \leq 0.5$ while driving BER to
$0.20$--$0.24$ at every SNR, against $0.005$ before.
The noise floor and the fading are the cover.
The learned half had to be rebuilt too, since a detector is only valid for the channel it was trained
on, and retrained on faded complex spectrograms it reaches $94.3\,\%$ accuracy at $91.2\,\%$ detection
and $2.3\,\%$ false alarms.
Composing the two per sample then reverses the lossless case, with the energy detector flagging nothing
the CNN misses across the entire sweep.
On a lossless channel the power meter subsumed the CNN, on a realistic one the CNN subsumes the power
meter, and that is where the expensive detector earns its keep.

\begin{figure}[!t]
\centering
% \includegraphics[width=\columnwidth]{fig/matched_detectability_crop.pdf}
\caption{Attainable BER against the detectability budget, blind versus channel-aware selection, at
5\,dB and 30\,dB. The two frontiers sit on top of each other.}
\label{fig:matched}
\end{figure}

The $70\,\%$ was measured at the same subcarrier count and the same power, which is a comparison at
matched attacker configuration.
Concentrating power where the link is faded raises BER and the received signature together, so a
channel-aware attacker is also a louder one and that comparison charged it nothing for the noise.
Holding fixed what the defender sees instead, we built the attainable frontier
$\mathrm{BER}^*(\varepsilon)$ -- the best BER over all configurations with suite
$\Pr[\text{flag}] \leq \varepsilon$ -- per strategy and per SNR.
A coarse grid gave gains swinging from $+5$ to $-57\,\%$, which is grid noise and not a measurement, so
we densified to nine powers by eleven subcarrier counts at 512 frames per cell.
The gains are then $-4.6$, $+1.5$, $-1.8$, $+6.2$ and $+0.0\,\%$ from 5 to 30\,dB.
Knowing the channel buys the genie nothing a blind attacker cannot get by spending the same
detectability, so channel-aware selection is not a lever, and learning to approximate the genie was
never going to be one either.

The same sweep costs us one more number, and this is the correction that matters most here.
The stealthy region above was read at $\Pr[\text{flag}] \leq 0.5$, and half of every frame is not
stealth -- that attacker is identified within a few frames, which contradicts the persistent BER floor
that motivated the step.
Read at the budget the defender sets, its own clean false-alarm rate ($\approx 0.12$ at 5\,dB falling to
$0.03$ at 30\,dB), no effective and unflagged attacker exists at 5--15\,dB, and at 20 and 30\,dB the best
attainable BER is $0.011$ and $0.004$, three and eleven times the clean floor.
The advantage is real and it is one to two orders of magnitude smaller than we reported.

Both retractions have the same cause: a comparison that held fixed something other than what the
defender sees.
The characterisation results are untouched by either, because none of them is defined against a stealth
threshold.
